\section{Finite-scale conventions and cited inputs}\label{sec:preliminaries}

\subsection{Directions, line families, and shadings}
Every ambient dimension $M$ is a fixed integer. Implicit constants may depend on $M$ and on a fixed spatial bound $R$,
but $M$ must not vary with the discretization scale.
Let $0<\delta<1/2$, and write
\[
  \Ldelta=1+\log(1/\delta).
\]
For unoriented unit directions $u,v$, use the projective angle
$\ang(u,v)=\arccos|\langle u,v\rangle|\in[0,\pi/2]$.
A set of directions is $\delta$-separated if every two distinct directions have projective angle at least $\delta$.

Let $\T$ be a finite labelled family of unit-scale line segments in a fixed ball $B(0,R)\subset\R^M$,
or of their neighbourhoods of width a fixed multiple of $\delta$.
A shading $Y$ assigns to each original line $T$ a collection of $\delta$-cells intersecting it or lying within distance $O(\delta)$ of it.
Write
\[
 m=\card{\T},\qquad E=\bigcup_{T\in\T}Y(T),\qquad
 S=\sum_{T\in\T}\card{Y(T)}.
\]
Here $\card{E}$ denotes the number of cells, not their Lebesgue volume.
Each original line has a fixed upper bound on its length, so
\begin{equation}\label{eq:single-line-upper}
  \card{Y(T)}\le C_{M,R}\delta^{-1}.
\end{equation}
Moving each cell to the nearest parameter point on its original line and removing multiplicities in parameter intervals of length $\delta$
loses only a fixed factor. We may therefore count separated marks on the original lines instead.
Throughout, cells and marks are interchanged only up to this fixed factor, without changing the small-power exponents.

For a real number $s\ge3$ and $A\ge1$, the required directional hypothesis is
\begin{equation}\label{eq:angular-cap}
 \card{\{T\in\T:\ang(v_T,u)\le\theta\}}
 \le A\left(\frac\theta\delta\right)^s,
 \qquad \delta\le\theta\le1,
\end{equation}
for every direction centre $u$.
This is only an upper bound; no pointwise Ahlfors--David lower bound, convex Wolff condition, or sticky hypothesis is imposed.
An average shading density $\lambda$ means
\begin{equation}\label{eq:mean-density}
 S\ge\lambda m\delta^{-1},\qquad 0<\lambda\le1.
\end{equation}
An average lower bound cannot directly replace a lower bound for each line.

\subsection{The four-dimensional finite directional-mass input}

\begin{inputthm}[Four-dimensional finite directional-mass estimate, archive DIR-003]\label{input:D4}
Fix a bounded four-dimensional graph chart. For every $\varepsilon_D>0$, there exist
$\eta_D>0$, $\delta_D>0$, and $C_D$ such that the following holds whenever $0<\delta<\delta_D$.
Let
\[
 \ell_i(t)=(t,b_i+t v_i),\qquad 0\le t\le1,\qquad
 |b_i|+|v_i|\le R_D,
\]
be a finite labelled line family, and let $J_i$ be its respective marks on the $\delta$ time grid, with
$\card{J_i}\ge\delta^{-1+\eta_D}$.
Given arbitrary positive weights $\omega_i$, suppose that every half-open slope cell $Q$ of side length $\delta$
satisfies $\sum_{v_i\in Q}\omega_i\le\alpha$.
Then the union $E_D$ of the cells visited by all the marks satisfies
\begin{equation}\label{eq:D4}
 \sum_i\omega_i\card{J_i}
 \le C_D\alpha\delta^{-\varepsilon_D}\card{E_D}.
\end{equation}
Repeated labels and arbitrary ranges of positive weights are allowed; the constants do not depend on the smallest weight or the number of labels.
Fixed-factor enlargements of the grid and visited cells are allowed, with corresponding changes in the constants.
\end{inputthm}

This input follows from a finite-scale consequence of the four-dimensional proof and a directional amplification;
see\cite{ArchiveSource,ArchiveDIR}.
It is not deduced directly from the set-dimension statement of\cite[Theorem 1.1]{OAI4} alone.
We use\eqref{eq:D4} with unit weights and separated output directions.
In particular, the input retains the order $\forall\varepsilon_D\,\exists\eta_D$;
no linear lower bound for $\eta_D$ in terms of $\varepsilon_D$ is assumed.

\subsection{Katz--Tao notation for estimates}

Set $N=\delta^{-1}$. We use the discrete-line, shading, and small-power-loss conventions of\cite[Definitions 5.1, 5.2, 6.1]{KT}.
For real $n,d>0$, $\Kest(n,d)$ means that, in every fixed ambient space $\R^M$,
a saturated $(n-1)$-collection of lines with saturated shading visits at least $N^{d-o(1)}$ cells.
The directional family satisfies the angular upper bound $N^{o(1)}(N\theta)^{n-1}$
for all centres and all $N^{-1}\le\theta\le1$, and has size $N^{n-1+o(1)}$.
Saturation of the shading means that its total mass is $N^{1-o(1)}\card{\T}$; saturation on each line is not assumed in advance.
The parameter $n$ is not the ambient dimension $M$ and need not be an integer.

Here $o(1)$ denotes only the arbitrarily small power losses allowed in the cited paper.
In an application, one first fixes the target output loss, then chooses a sufficiently small input tolerance, and finally takes $N$ sufficiently large.
Section~\ref{sec:seeds} gives an explicit $\forall\varepsilon\,\exists\eta$ sufficient form for this use.
We invoke only the strictly subcritical $\Kest(n,4-\zeta)$, without asserting a uniform linear modulus relating the critical input and output losses.

\begin{inputthm}[Katz--Tao recurrence, Theorem 6.2 of\cite{KT}]\label{input:KT}
Let $0<d<n$. If $\Kest(n,d)$ and $\Kest(d+1,d')$ hold, then
\begin{equation}\label{eq:KT}
  \Kest\!\left(n,\frac{2n+1+d'}4\right)
\end{equation}
holds. The input allows real parameters and every fixed ambient space.
\end{inputthm}

The second application in this recurrence may take place in $\R^{M+1}$ and requires the unsaturated version in equation (40) of the cited paper:
\begin{equation}\label{eq:KT-unsaturated}
  \card{E}\ge
  \frac{\card{\T}}{N^{n-1}}N^{d-o(1)}.
\end{equation}
Thus a saturated conclusion in just one fixed ambient space is insufficient.
Below we retain the original number-of-lines factor directly, establishing this interface;
the full proof of the recurrence remains a cited input from\cite{KT}.
